% `tex.web` §811: every cell in a finished alignment row takes the *row's*
% height and depth, not its own.
%
%     @<Make the unset node |r| into an |hlist_node| of width |w|…@>=
%     begin height(r):=height(q); depth(r):=depth(q);
%
% A cell is a window onto the row rather than a box of its own size. That is
% what makes a `\vrule` with a running height fill a table row however short
% the entry holding it: the rule stays running — `\rule(*+*)` in a display, in
% both engines — and the *cell* it is measured against grows to the row.
%
% We left each cell at its natural size, so a cell holding only a rule was 0pt
% tall and the rule had nothing to run to. Found by diffing trip.tex's sixth
% shipped page, whose first `\hbox` is 22.0 tall in Knuth's transcript and was
% 0.0 here.
\catcode`\{=1 \catcode`\}=2 \catcode`\#=6 \catcode`\&=4
\font\r=cmr10 \r \hsize=200pt
\showboxdepth=3 \showboxbreadth=9 \tracingonline=0

% a rule cell beside a tall cell: both cells are the row's height
\setbox0\vbox{\halign{#&#\cr \vrule & \vbox to 22pt{}\cr}}
\showbox0

% depth travels with it, not just height — the lowered box gives the row one
\setbox1\vbox{\halign{#&#\cr \vrule & \lower3pt\vbox to 22pt{}\cr}}
\showbox1

% a row whose entries differ: the short one is raised to the tall one
\setbox2\vbox{\halign{#&#\cr \vbox to 5pt{} & \vbox to 30pt{}\cr}}
\showbox2

% ...and two rows do not borrow from each other
\setbox3\vbox{\halign{#\cr \vbox to 5pt{}\cr \vbox to 30pt{}\cr}}
\showbox3

% §810's vertical branch is the transpose: a `\valign` cell takes the
% *column's* width, and a cell holding a running `\hrule` was 0pt wide here
% with the rule having nothing to run to.
\setbox4\hbox{\valign{#&#\cr \hrule & \hbox to 13pt{}\cr}}
\showbox4

% ...and two columns do not borrow from each other either
\setbox5\hbox{\valign{#\cr \hbox to 5pt{}\cr \hbox to 30pt{}\cr}}
\showbox5

% §768: a `\valign` entry is a *vertical* list, so text in it starts a
% paragraph — and the entry's end is what breaks that paragraph into lines,
% exactly as a `\vbox`'s closing brace is. Our alignment path closed the cell
% without ending the paragraph, so a `\valign` cell containing text closed with
% it still open and the content went nowhere: an empty cell.
%
% An `\halign` entry is horizontal and has no paragraph to end, which is why
% this is asked of the vertical case only — the `\halign` cases above cover
% that side.
\setbox6\hbox{\valign{#\cr A\cr}}
\showbox6
\setbox7\hbox{\valign{#&#\cr A\cr}}
\showbox7
\end
